\section{Related work}
\label{sec:related}

\subsection{Engagement and attention recognition}
Automatic engagement recognition grew out of facial-expression analysis~\cite{whitehill2014faces} and has been reviewed repeatedly~\cite{karimah2022automatic,booth2023engagement,qarbal2025student}. Webcam datasets shaped the field: DAiSEE~\cite{gupta2016daisee} provides four-level labels for short clips, EngageNet~\cite{singh2023engagenet} provides about 11{,}300 ten-second clips of 127 participants, and the EmotiW challenges made engagement prediction a shared task~\cite{dhall2018emotiw,dhall2020emotiw}. On the official DAiSEE test set, reported four-class accuracy has risen from 57.9\% for the baseline of the dataset paper~\cite{gupta2016daisee} to about 67\% in peer-reviewed work~\cite{liao2021dfstn,abedi2021resnettcn,selim2022students,abedi2024affectordinal,abedi2025supervised}, and a recent preprint reports 73\%~\cite{gothwal2025vibednet}. Most of these results are single accuracies on one split, without seed variation or a chance-corrected metric, and ordinal structure is exploited only occasionally~\cite{abedi2024affectordinal,abedi2025supervised}. Much higher accuracies have been reported on DAiSEE after class rebalancing under other splits~\cite{santoni2024automatic}, and on newly collected datasets~\cite{selim2022students}; such figures are not comparable with the official split. Landmark trajectories alone are competitive on EngageNet~\cite{yulius2026student}, and zero-shot vision--language models stay near chance-level agreement on DAiSEE while their accuracy sits at the majority rate~\cite{goyal2026zeroshot}.

Work in physical classrooms is rarer. S\"umer et al.~\cite{sumer2023multimodal} recorded secondary-school classes and obtained modest person-independent performance from head-pose and facial-expression features, with a fusion on par with the best modality and a gain from a short personal calibration. Goldberg et al.~\cite{goldberg2021attentive} validated manual ratings of visible engagement in university lectures and found that a feature describing the synchrony between students improves their estimation. Wearable and environmental sensors have been used for self-reported engagement in a high school~\cite{gao2020ngage}, and classroom image and clip datasets have been released for behavioural categories and group engagement~\cite{xu2023icapd,lu2025video}. Mind-wandering detection from face video is only moderately above chance and degrades from the laboratory to lecture viewing~\cite{buhler2025lab}, and an independent evaluation of wearable affect detection in everyday life found performance close to zero~\cite{dmello2023affect}.

DIPSER~\cite{marquezcarpintero2025dipser} is, to our knowledge, the only available dataset with per-student video, landmarks, head pose and smartwatch signals recorded in in-person classes and rated by several annotators and by the students. Its authors studied vision--language models on its images~\cite{marquezcarpintero2025enhancing}; Liu et al.~\cite{liu2025attention} combined facial and wearable data; Fragkiadakis et al.~\cite{fragkiadakis2026fairness} found that a wearable-sensor model does not improve on the training mean; and our earlier work provided single-stream face baselines~\cite{makhkhas2026ordinal}. One recent study reports a macro-F1 of 0.77 on DIPSER with a large gain from fusion and from physiological signals~\cite{dan2026multimodal}. Its description of the data differs from the dataset description, so we could not relate its figures to ours; the split comparison of Section~\ref{sec:fc-protocol} shows how strongly such figures depend on the protocol.

\subsection{Multimodal fusion and its failure modes}
Fusion strategies range from early fusion by feature concatenation to late fusion of decisions~\cite{baltrusaitis2019multimodal}; gated units learn input-dependent weights~\cite{arevalo2017gated}, and modality dropout improves robustness to missing streams~\cite{neverova2016moddrop}. Several studies show that joint training does not guarantee a gain. Multimodal networks can underperform their best unimodal counterpart because they overfit more and because modalities overfit and generalise at different rates~\cite{wang2020makes}; a dominant modality suppresses the gradients of the others~\cite{peng2022balanced}; networks learn modalities greedily~\cite{wu2022characterizing}, which has been analysed theoretically as modality competition~\cite{huang2022modality}; and transformers degrade when a modality is absent~\cite{ma2022robust}. Hessel and Lee~\cite{hessel2020does} show that apparent multimodal gains are often additive and recommend always reporting unimodal baselines. In education, reviews of data fusion in multimodal learning analytics describe many fused systems~\cite{chango2022datafusionmmla}, and fusion with audio has been reported to help for engagement in online presentation classes~\cite{wu2024cmose}. CORE-Net~\cite{tran2026ordinal} recognises ordinal engagement from several modalities in collaborative learning. The present study differs from these works in that the fused and single-stream models are compared under one protocol with shared front-ends, folds and seeds.

\subsection{Ordinal objectives and multi-rater labels}
CORAL obtains rank-consistent predictions for a $K$-level problem from $K-1$ binary tasks with shared weights~\cite{cao2020coral}, CORN relaxes the weight sharing~\cite{shi2023deep}, and CORE-Net combines a cross-entropy, a squared earth mover's distance and a neighbour penalty~\cite{tran2026ordinal}. Labels from several raters can be voted, aggregated by expectation--maximisation~\cite{dawid1979em} or retained as distributions~\cite{geng2016ldl}; learning from disagreement has been surveyed in~\cite{uma2021learning}. Engagement labels are a known difficulty: datasets differ in definition, source and temporal resolution~\cite{khan2022inconsistencies}, observer ratings dominate~\cite{karimah2022automatic}, and a recent audit of DAiSEE and EngageNet finds that the labels of adjacent clips agree weakly on DAiSEE and that about half of the label variance of EngageNet lies between participants~\cite{le2026auditing}. That audit is the closest work to ours on the two webcam datasets. It does not address classroom data, sensors, fusion, or the comparison of splits that hold out students and classes.

\subsection{Evaluation protocols and leakage}
Dependence between training and test data is a general cause of optimistic results~\cite{kapoor2023leakage}. With repeated measures, record-wise splits inflate scores~\cite{saeb2017need} because the model can identify the subject~\cite{chaibubneto2019detecting}; the same has been shown for segments of physiological recordings~\cite{brookshire2024data,pandilova2026participant}, for sliding windows in activity recognition~\cite{hammerla2015lets,tello2023too}, where the gap between ten-fold and leave-one-subject-out accuracy can exceed ten points~\cite{gholamiangonabadi2020deep}, and for micro-expression recognition~\cite{varanka2024data}. Slow temporal structure can also be mistaken for the target~\cite{li2020perils}, and structured data call for blocked cross-validation with the block chosen according to the intended use~\cite{roberts2017cross}. Selection on the scored data adds a bias of comparable size to differences between algorithms~\cite{cawley2010overfitting}. Trivial references matter as well: personal baselines are often omitted~\cite{demasi2017meaningless}, the no-change forecast is a required benchmark for series~\cite{hewamalage2023forecast}, and for self-rated attention in video learning, carrying the last rating forward gives a lower forecasting error than sensor models~\cite{leng2026edugage}. In multimodal learning analytics, reviews call for full disclosure of evaluation practice and for caution about what sensor data can support~\cite{worsley2021new,yan2022scalability,cukurova2020promise}, noting that sensors capture behaviour whereas cognitive states are latent and reach the model only through annotation~\cite{dimitri2018signals}. We add to this literature a measurement, on one classroom dataset and with the model held fixed, of the three levels of dependence that such data contain: the window, the student and the class.

\subsection{Discrete orthogonal moments}
Discrete Tchebichef moments provide compact and numerically stable image descriptors without numerical approximation of continuous integrals~\cite{mukundan2001tchebichef}. They are used here as a fixed front-end, so that every model of the study is small and all streams are described by the same transform.
